\subsection*{第7章の索引用語}

\begin{description}
  \item[\term{ラベル}] 教師あり学習で入力へ対応づける正解または目標値です。
  \item[\term{L1正則化}] パラメータの絶対値和を損失へ加える正則化です。
  \item[\term{線形分離}] 一つの超平面によってクラスを分けられる性質です。
  \item[\term{ニューロン}] 重み付き和と活性化関数で入力を変換する計算単位です。
  \item[\term{ソフトマックス}] 複数の実数を合計1の正の値へ変換する関数です。
  \item[\term{位置情報}] 系列中で各トークンが占める順序や距離の情報です。
  \item[\term{Decoder-only}] Decoderの自己回帰処理を中心に構成するTransformerです。
  \item[\term{BPE}] Byte Pair Encodingの略で、頻出する単位を反復的にまとめるトークン化方式です。
  \item[\term{ファインチューニング}] 事前学習済みモデルを目的別データで追加学習することです。
  \item[\term{文章生成}] 次トークンの選択を繰り返して系列を作る処理です。
  \item[\term{偏り}] データや評価条件の偏在によって結果が群や条件ごとに系統的に変わることです。
  \item[\term{汎化}] 学習時に見ていない入力へも有用な予測を返す性質です。
\end{description}
